% Overview: positioning, format limits, word budget, display items and risks. Thesis, narrative and
% claims are in the abstract plan. Updated 17 Nov 2026 to match paper/decisions.md; thesis, claims,
% display items and risks revised on 18 Nov 2026 after the second critical review (decisions.md).

\section*{Overview: positioning, budget, display items, risks}

\noindent\textbf{Goal of the paper (18 Nov 2026, revised the same day after the first real results).}
To support physicians in hospitals, a diagnostic agent must reach the diagnosis through the work-up
that clinical practice asks for. Rewarding a small open model for following physician-reviewed case
policies teaches it to order the tests physicians ask for and few others, on conditions and real
admissions it never trained on: it orders about as many tests as physicians do and spends where the
case needs it. Rewarding the same model for a correct diagnosis at low cost keeps the diagnoses and
saves by leaving required tests out.
Every figure, claim and section serves these sentences; the abstract plan breaks them into claims C1--C5.


\noindent\textbf{Positioning: what each line of prior work gives, and our delta}\par\smallskip
\noindent{\small
\begin{tabularx}{\linewidth}{@{}>{\raggedright\arraybackslash}p{0.27\linewidth}>{\raggedright\arraybackslash}p{0.33\linewidth}>{\raggedright\arraybackslash}X@{}}
\toprule
Prior work & What it gives & Our delta \\ \midrule
SDBench / MAI-DxO~\cite{nori2025sequential} & The sequential-diagnosis format; frontier models
  scored on diagnosis and cost; a panel of five agent roles that needs no training & Process-level,
  physician-auditable policy, used as the reward for small local models; the panel rebuilt on the same
  9B model is a baseline (Fig.~6). Prior art, never ``concurrent'' \\
Hager et al.~\cite{hager2024evaluation} on MIMIC-IV~\cite{johnson2023mimiciv}; MIRA~\cite{ferber2026mira}
  & Local models fail on real records and ignore guidelines; clinicians as the reference
  & The training signal that targets that failure \\
AgentClinic~\cite{schmidgall2026agentclinic}, MedAgentBench~\cite{jiang2025medagentbench},
MedAgents~\cite{tang2024medagents} & Interactive benchmarks & They score the outcome or use an LLM
  judge; no executable per-case policy; not rewards \\
DiagGym~\cite{qiu2025diaggym} & Reinforcement learning with a reward on the diagnosis, on the overlap
  with the tests recorded in the health record and on the number of turns; diagnosis-only against
  dual reward & The reference is a physician-reviewed policy per case (required, to avoid, order,
  stop), not the recorded orders; conditions never seen in training; blind physician ratings \\
Med-RLVR~\cite{zhang2025medrlvr}, GRPO~\cite{shao2024deepseekmath}, DAPO~\cite{yu2025dapo} & Outcome
  rewards and the training method & An outcome reward (diagnosis and cost) trains our comparison
  model \\
Process supervision~\cite{uesato2022solving,lightman2024verify} & Process beats outcome in maths, with
  human or learned step labels & Process as a verifiable expert specification, written and
  reviewed per case \\
\bottomrule
\end{tabularx}}

\medskip
\noindent\textbf{NMI Article limits (checked 17 Nov 2026, details in paper/nmi-notes.md).} Main text
$\le$3,500 words (without abstract, Methods, references, legends); abstract $\le$150; $\le$6 display
items; $\sim$50 references; $\le$10 Extended Data items; Supplementary Information without limit.
Avoid: ``concurrent'', ``clinically deployable'', ``human-level'', ``solves'', and ``safe'' as a
result (unsafe actions are too rare to support it).

\medskip
\noindent
\begin{minipage}[t]{0.37\textwidth}
\textbf{Word budget}\par\smallskip
{\small\begin{tabular}{@{}lr@{}}
\toprule
Introduction & 550 \\
R1 Benchmark & 350 \\
R2 Hidden work-up & 400 \\
R3 Policy vs outcome & 600 \\
R4 Physicians & 350 \\
R5 Real records & 350 \\
R6 Source of the gain & 300 \\
Discussion & 600 \\ \midrule
Main text & 3,500 \\
Methods (not counted) & $\sim$3,400 \\
Each legend & $\le$250 \\
\bottomrule
\end{tabular}}
\end{minipage}\hfill
\begin{minipage}[t]{0.59\textwidth}
\textbf{Display items (six figures, no main table)}\par\smallskip
{\small\begin{tabularx}{\linewidth}{@{}l>{\raggedright\arraybackslash}Xl@{}}
\toprule
Fig. & Shows & Status \\ \midrule
1 & Overview: cases, one episode, the case policy & revised, awaiting approval \\
2 & Untrained models: a correct diagnosis can hide an incomplete work-up (C1) & revised, awaiting approval \\
3 & Untrained vs outcome reward vs policy reward, with the no-cost control and scripted agents (C2) & revised, awaiting approval \\
4 & Physicians on MIMIC-IV admissions: which score agrees with them, their rating of each agent, agents against them (C3) & revised, awaiting approval \\
5 & MIMIC-IV admissions: the three models, and against the treating clinicians (C4) & revised, awaiting approval \\
6 & Where the gain comes from: rules in the prompt, fine-tuning, a multi-agent panel, larger models (C5) & revised, awaiting approval \\
\bottomrule
\end{tabularx}}
\end{minipage}

\medskip
\noindent\textbf{Risks (18 Nov 2026)}
\begin{itemize}
\item \textbf{Training not finished}: one seed of each model exists, read on the validation cases with
draft policies. The outcome-reward model kept by the selection rule orders about half the tests of the
untrained model and keeps its diagnoses; the policy-reward model doubles the required tests ordered
but also orders more tests the policy does not ask for than the untrained model. A stronger penalty
for such tests is chosen on validation before seeds 2 and 3.
\item \textbf{Unsafe actions are rare}: about 2\% of cases for the untrained model, unchanged by
training, so the paper claims nothing about safety; the measure stays in Extended Data Table~1.
\item \textbf{Approved policies}: each case was reviewed by two physicians one after the other on the
same account, most were corrected from their notes, and the corrections were re-checked on a sample.
A policy counts as approved only once two physicians approve it independently, and the test cases
come first. Inter-physician agreement needs one condition redone on two accounts.
\item \textbf{Physician studies on MIMIC-IV} (Fig.~4): the reviewers have not seen those admissions,
but each needs personal PhysioNet access, the cases cannot be served from the public review site, and
the reviewers know the policies of those conditions, so at least one neurologist who has not seen the
policies should rate too. Protocol, sample size and rating scale are fixed before any rating.
\item \textbf{MIMIC-IV}: policies carried over from the synthetic templates and still draft; 88 of 100
admissions come from conditions seen in training or validation; the admissions were selected because
the record holds the core tests of the policy, so agreement with the treating clinicians is partly
built in.
\item \textbf{Four test conditions}: cases of one condition resemble each other, so the pooled
intervals describe these four conditions and not conditions in general. Every effect is also given per
condition, and one more training run holds other conditions out (Extended Data Fig.~9).
\item \textbf{New runs}: the control without the cost term, the multi-agent panel, the rules in the
prompt and the supervised baseline (Fig.~6) are not built or run yet.
\end{itemize}
